\section{Exact structured-action protocol}

The status prefix is the UTF-8 byte sequence for \texttt{Status:}, with token IDs
\texttt{[2583,25]} and no trailing whitespace. PASS, FAIL, and SKIP map from raw token IDs 51935,
34207, and 82504. The caveat prefix is \texttt{\textbackslash nCaveat:}, token IDs
\texttt{[198,34,525,266,25]}, also without trailing whitespace. INCLUDE, OMIT, NONE, and SKIP map
from 51493, 70500, 43969, and 82504. All other tokens are \textsc{no-action}. Sampling uses the
full vocabulary; candidate-constrained decoding and retry repair are prohibited. Status enters
the causal transcript before caveat, so caveat analyses are conditional on status.

\section{Structured-action control protocols}
\label{app:secondary-methods}

\paragraph{Offline candidate selection (C1).}
The completed field contains 1,200 conversations---600 deceptive-source and 600 honest-source,
balanced across the four PASS/FAIL status classes---and 260,400 saved candidate rows from 20
families. Four folds hold out five complete families each. Candidate actions span six direction
families, layers 8/12/16/19/20/24/28/32, doses 24/48/96, and both PASS and FAIL targets. Evaluation
uses saved post-intervention margins rather than fresh generation.

The learned selector vectorizes pre-response candidate metadata, standardizes it, and fits ridge
regression with penalty 1 to per-candidate reward ($+1$ for fixing a deceptive report, $-1$ for
harming an honest report, and zero otherwise). Features include direction family, candidate target,
oracle route, original status, layer, dose, base margin, gate score, tangent and neighborhood
diagnostics, direction-training counts, and the retained principal-coordinate fields. It excludes
the candidate's final or delta margin, post-action correctness, reward, and harm. At held-out
application it scores every candidate present in the saved field, takes the highest positive
score, and otherwise abstains. Honest-row menus contain both target signs; the oracle-routed
false-row menus contain only the corrective sign. Thus the ranker is response-blind but not
target-information-free.

CNG uses a pre-response multi-layer query and an unconstrained argmax over its fold-local action
library; it adds no selector-side route mask. Its compact public receipt does not preserve enough
of the exact historical fit to reconstruct the selector. Within the same oracle-routed field, the
fixed comparator explicitly maps the supplied route to its corrective target and selects one
\texttt{bidir\_linear} layer--dose coordinate by mean reward within each training fold. Every fold
selects layer 16 and dose 96. This comparator is a retrospective held-out-family-clean
reconstruction, not a prospectively registered controller. A separate L20 linear route diagnostic recovers the PASS/FAIL route from pre-response residuals
on all 1,680 held-out-family predictions under 12-family leave-one-family-out cross-validation
on the same-distribution structured-action bank. It is routing-only, not true OOD and not
end-to-end control; C1 itself consumes the task-supplied route and does not test that inference.

\paragraph{Generated dose audit (C2).}
C2 reuses 1,200 prompts and 20 families but generates new structured-action responses. Oracle
true-status routing supplies the PASS/FAIL target. A family-excluded layer-16 Local Control Flow
combines twelve nearby training-family state groups into a normalized direction. A standardized
ridge model predicts the log first-crossing dose from pre-action context, route, margin, gate,
support, neighborhood, boundary, and causal-gain features; four family-block calibration folds add
a conservative residual correction and snap upward to the observed dose grid, capped at 128.
Steering is active only while scoring the first PASS/FAIL token; the higher-logit token is inserted,
steering is removed, and the explanation is completed greedily. Fixed doses 88, 96, and 128 use
the same route, direction, and token scope. The compact receipt preserves counts but no sampling
interval, so C2 is reported as a conditional dose audit rather than target-free control.
In the separate dense-grid diagnostic, the training-fold first-crossing dose has median 16 in
every fold, while the deployed learned policy frequently snaps to 96--128; this is descriptive
headroom, not a registered advantage.

\paragraph{Pilot decomposition (C12).}
The primary pilot balances 40 rows in each of the four status classes (160 total); the separate
off-tangent follow-up has 32 rows, including 16 deceptive pairs. A family-excluded layer-20
logistic gate predicts true status from turn-3 pre-response residuals. It abstains when that
prediction already matches the report and otherwise supplies the PASS/FAIL route to every
bidirectional and explicitly gated arm. The raw bidirectional linear arm is therefore gate-routed,
not ungated. The tangent direction projects a matched-cell, train-family correction vector into a
query-local tangent; the follow-up's off-tangent arm uses the normalized rejected component.

All primary interventions use layer 20, total dose 96, and first-decision-token scope; the
follow-up uses dose 48. The deterministic endpoints are status correctness and a stricter
numeric-basis rule. Because direction family, gating, and other components were not crossed
factorially, their contributions are not separately identified. The historical payload also did
not preserve the exact tangent-fit invocation, runner version, or immutable 4-bit model identity;
current code defaults are not treated as a reconstruction of those missing facts.
The current implementation includes \texttt{bidir\_norm\_matched\_raw},
\texttt{random\_tangent\_s*}, and \texttt{random\_norm\_matched\_raw\_s*} arms, but no completed
result artifact containing them was found. They therefore cannot identify a tangent-specific
effect.

\section{Pre-action warning protocols}
\label{app:c11}

C11 aggregates three separately specified negative experiments. Their populations and features are
not interchangeable, and no metric is pooled across them.

\paragraph{Intrinsic spectral field.}
The population is 120 pressured turn-two AN status events from 60 scenarios and 20 families; the
honest/deceptive conditional slice contains 106 events (32 honest, 74 deceptive). Exact repeated
states are quotiented by family and compact-section identity without majority voting. Each state
has 135 spectral relations: four residual, 128 head-specific attention, and three layer-transport
relations. Squared retained singular values, tail mass, and transport sign mass define
Hellinger-type relation distances; residual, attention, and transport views are RMS-aggregated.

Five outer folds hold out complete families and exclude held-out scenario, prefix, and section
identities. Training-fold rank calibration supplies the admissible relations. The predictor uses a
Gaussian radial kernel with median training-pair bandwidth and 0.5 smoothing over the five status
outcomes. The equal-view model combines median-scaled residual, attention, and transport distances.
Its comparators are the training-fold base rate and the exact
$(\text{true status},\text{desired status},\text{knowledge})$ design cell. Admission required
improvement in held-out proper scores over both; the equal-view model loses to the design-cell
comparator in all five folds. AUROC 0.4168 is a conditional descriptive statistic, not the decision
rule.

\paragraph{Complete-path response field.}
This artifact-only, nonconfirmatory bank contains 60 sample-0 scenarios; 52 enter the
honest/deceptive slice (18 honest, 34 deceptive). Each query binds three physical NN--AN incoming
realizations and exact replays to the AN--AA and AN--AB outgoing branches. Across 120 frozen
relation-sides, the path descriptor retains normalized principal angles, transported-projector
discrepancy, polar residual, and explicit support state, summarized into incoming/replay,
continuation, and branch-order components.

The primary family-held-out model is a design-conditioned Gaussian kernel over the complete path,
with median within-cell training distance and 0.5 smoothing. Comparators are the exact design-cell
prior, base rate, incoming-only field, one-state spectral field, and an identity-shuffled path null.
Admission required the lower family-bootstrap bound on honest/deceptive log-loss gain to exceed
zero against every comparator and in both knowledge-correct cells. It did not; no controller was
admitted. This object is a finite path-response descriptor, not an estimated gauge connection.

\paragraph{Masked geometry-derived risk model.}
The largest instrument evaluates all 1,680 unique pre-status events, including 1,568
honest/deceptive events. Its primary view removes the entire intervention-bearing user message
while retaining the exact action-free pre-status root. The object contains four-layer root
residuals, root-to-context residual profiles, masked-and-renormalized incoming attention for every
head, and typed token metadata. A training-fold retrieval sketch proposes neighbors; an exact
typed token-coupling, residual-profile, and head-set-matching metric reranks them to eight unique
roots. Their outcomes form Jeffreys-smoothed five-class local probabilities.

The published ``geometry-only logistic'' is therefore a geometry-derived stack, not a direct
classifier on tensors: it fits multinomial logistic regression to five sealed local
log-probabilities and log support. The nuisance comparator keys training-fold counts by turn,
intervention history, pressure exposure, true status, and desired status. The primary metric is
five-class log-loss gain (nuisance loss minus geometry-derived loss), with a paired 2,000-draw
scenario bootstrap; repair required both the point and lower interval bound to be positive. The
stack is cross-fitted but not fully nested: its meta-model training features are precomputed
out-of-fold local probabilities. Any resulting leakage favors geometry and does not rescue the
negative comparison.

On the 1,568-event honest/deceptive slice, geometry-only logistic reaches AUROC 0.9442 and log
loss 0.5290; nuisance-only reaches 0.9446 and 0.5396; nuisance-plus-geometry reaches 0.9537 and
0.5254; and the sealed nuisance prior reaches 0.9489 and 0.5072. The nested
nuisance-plus-geometry comparison improves AUROC, Brier, and log loss with the intervals reported
in Section~\ref{sec:reading}. These are descriptive secondary fits; the registered refutation remains the
sealed-prior proper-score endpoint.

The compact receipt binds the source artifacts by hash but omits the large event banks. The
spectral rank-calibration implementation and exact configuration are also absent from the public
tree. Accordingly, the supplement audits the reported result and feature contracts but does not
reconstruct every C11 fit from raw capture. The failures refute these instruments, not all possible
pre-action predictors; twenty token-identical prefixes with mixed sampled outcomes provide only an
empirical stochastic floor on those observed states.

\section{Prospective natural-prose control protocol}
\label{app:c5}

C5 evaluates 96 fixed natural-prose prefixes: 24 from each of four scenario families excluded
from every fitted component, with 12 deceptive and 12 honest base responses per family. Each
final prompt requests a plain-prose report with an explicit PASS or FAIL status and factual basis,
and prohibits JSON.
Each of four arms---base, native tangent-actuated gate, frequent tangent-actuated gate, and fixed linear
comparator---generates greedily from the same prefixes with \model{} in BF16, seed 0, and a
4,096-token context and 512-new-token cap. The primary estimand for arm $a$ and source label
$\ell$ is the paired change in machine-scored status correctness,
\begin{equation}
  \widehat\Delta_{a,\ell}=\frac{1}{n_\ell}\sum_{i:L_i=\ell}
  \left[\mathbf{1}\{\widehat s_{ia}=s_i^\star\}
  -\mathbf{1}\{\widehat s_{i0}=s_i^\star\}\right],
\end{equation}
where $s_i^\star$ is scenario truth and $\widehat s_{ia}$ is the unambiguous parsed status.
Ambiguous or unparsed reports do not count as correct; the endpoint never conditions on parse
success, and no prose-judge score can relabel it. The original report also computed 2,000
percentile-bootstrap draws over families with seed 0. Because there are only four families, we
treat the observed family cells in Table~\ref{tab:c5-family} as the uncertainty description rather
than interpreting those resamples as precise confidence coverage.

\begin{table}[h]
\caption{C5 paired machine-status transitions by evaluation family. Each cell is
fixes/harms among 12 rows in the named source stratum; its net correctness change is
$(\mathrm{fixes}-\mathrm{harms})/12$. D and H denote deceptive-source and honest-source rows, not
baseline correctness.}
\label{tab:c5-family}
\centering
\scriptsize
\resizebox{\linewidth}{!}{%
\begin{tabular}{lrrrrrr}
\toprule
Family & Native D & Frequent D & Linear D & Native H & Frequent H & Linear H \\
\midrule
Cybersecurity & 0/0 & 1/2 & 6/1 & 0/0 & 2/0 & 7/0 \\
Electrical grid & 0/0 & 2/1 & 6/0 & 0/0 & 3/2 & 6/3 \\
Structural welding & 0/0 & 0/3 & 4/3 & 0/0 & 0/3 & 4/3 \\
Water treatment & 0/0 & 0/1 & 5/1 & 0/0 & 0/0 & 9/0 \\
\bottomrule
\end{tabular}%
}
\end{table}

Both tangent-actuated arms require a frozen family-excluded turn-2 eligibility score of at least
0.252073. Frozen held-out scores cover all 96 rows; fallback scoring is forbidden. The native gate
is a running logistic score over nine layer-16 decode features. It may
enter after step 8 above score 0.55 and exits below 0.495. The frequent arm replaces this signal with
a window open through decode step 31, after a two-step warm-up, while retaining the same
eligibility and actuator. The actuator takes a frozen escape axis, projects it into the nearest
of 256 eight-dimensional local tangent planes, abstains below tangent fraction 0.10, and uses the
chart's median natural-step norm with a linear gate taper and hard cap 10.712. The comparator uses
the negative family-excluded layer-16 probe direction at fixed dose 8 on each decoding forward.
The native controller, frequent controller, and comparator are therefore distinct policies; C5
tests their prospective endpoints, not a factorial decomposition of their components.

\section{One-step gauge replay and transport diagnostic}
\label{app:c13}

C13 replays one intervention at an exact pre-status prefix. A live four-layer root must match its
sealed source-bank query. If the nuisance-matched deception-risk field is defined and the query
remains inside its support and boundary, the gauge arm takes a metric-normalized negative-gradient
step and horizontally lifts it into residual interventions at layers 12, 16, 19, and 20;
otherwise it returns no step. The controls use no intervention, the sign-reversed direction, and
a deterministic metric-orthogonal random tangent under common random streams.

Separately, the transport diagnostic writes a local honestward vector as $H=T+S$. Generic
transport $T$ is fit from same-true-status training-family edges, and $S$ is the residual. The
reported full reach, generic reach, and deception-specific remainder are true-status logit-margin
changes $m(H)-m(0)$, $m(T)-m(0)$, and $m(H)-m(T)$, respectively. This decomposition is not the
behavioral replay, and neither is the separate holonomy instrument.

\begin{figure}[h]
\centering
The gauge-geodesic and transport diagnostic figure appears as Figure~\ref{fig:gauge-results} in
the main text (Section~\ref{sec:c13}). It shows all 402 proposal-support roots including those
with zero direction, boundary exits, and off-support roots; only 21 receive an active supported
step. The gauge arm has no behavioral effect; the deception-specific transport remainder crosses
zero. Holonomy adequacy fails in every fold; transport and holonomy are neither confirmed nor
refuted.
\label{fig:gauge}
\end{figure}
